\documentclass[12pt]{article}
\usepackage[T1]{fontenc}
\usepackage[utf8]{inputenc}
\usepackage{lmodern}
\usepackage{textcomp}
\usepackage{enumitem}
\usepackage{geometry}
\geometry{margin=1in}
\begin{document}
\begin{center}
Answer Key - Business Administration - Graduate Level Assessment 21 \
\small (Answers are ordered exactly as in the question paper.)
\end{center}

\section*{Multiple Choice}
\begin{enumerate}[label=\arabic*.]
\item \textbf{Answer:} A
\\ \textit{Options:} A) 14\%; B) 17.5\%; C) 10\%; D) 15\%; E) 19.8\%
\\ \textit{Note:} The annual interest rate is calculated by determining the total cost of the radio on the installment plan, which is $5.00 down plus $3.10 per month for 12 months, totaling $43.20, and comparing this to the original price of $39.90, resulting in an effective interest rate of approximately 14\%.
\item \textbf{Answer:} A
\\ \textit{Options:} A) 3.5; B) 1.5; C) 1; D) 6; E) 5
\\ \textit{Note:} The current ratio is calculated by dividing current assets by current liabilities, so with current assets of $150,000 and current liabilities of $50,000, the current ratio is 150,000 / 50,000 = 3.0.
\item \textbf{Answer:} D
\\ \textit{Options:} A) Corporate brands.; B) Manufacturer brand.; C) Sub-branding.; D) Family branding.; E) Mixed branding.
\\ \textit{Note:} Family branding refers to a marketing strategy where a company uses its brand name for multiple related products, enhancing recognition and trust among consumers, as seen with companies like Microsoft, Heinz, and Kellogg's.
\item \textbf{Answer:} A
\\ \textit{Options:} A) 19.00; B) 20.35; C) 25.40; D) 23.50; E) 24.75
\\ \textit{Note:} The intrinsic value of a preferred share can be calculated using the formula for valuing a perpetuity, which states that the intrinsic value equals the annual dividend divided by the required rate of return; thus, $1.90 divided by 0.09 equals $19.00.
\item \textbf{Answer:} A
\\ \textit{Options:} A) Multi-product branding.; B) Sub-branding.; C) Mixed branding.; D) Family branding.; E) Corporate brands.
\\ \textit{Note:} The correct answer is "multi-product branding" because this strategy involves branding each product independently, allowing for distinct brand identities that cater to different markets or consumer preferences.
\end{enumerate}

\section*{True / False}
\begin{enumerate}[label=\arabic*.]
\setcounter{enumi}{5}
\item \textbf{Answer:} False
\\ \textit{Options:} A) True; B) False
\\ \textit{Note:} The statement is false because the cost per thousand for a full-page ad in Punk Magazine is typically higher than \$15, reflecting the magazine's niche market and lower circulation compared to larger mainstream publications.
\item \textbf{Answer:} True
\\ \textit{Options:} A) True; B) False
\\ \textit{Note:} Understanding customer willingness to pay is crucial for effective sales pricing because it allows businesses to align their pricing strategies with perceived value, maximizing revenue potential rather than merely covering production costs.
\item \textbf{Answer:} False
\\ \textit{Options:} A) True; B) False
\\ \textit{Note:} Carroll's framework for corporate social responsibility includes Economic, Legal, Ethical, and Philanthropic categories, not Political or Environmental.
\item \textbf{Answer:} False
\\ \textit{Options:} A) True; B) False
\\ \textit{Note:} Businesses engage in social accounting primarily to fulfill regulatory requirements, enhance their corporate reputation, and demonstrate accountability to stakeholders, rather than solely focusing on satisfying customer needs and improving relations.
\item \textbf{Answer:} True
\\ \textit{Options:} A) True; B) False
\\ \textit{Note:} The price-to-earnings (P/E) ratio is calculated by dividing the market price per share by the earnings per share, so for Company A, the P/E ratio is $150 ÷ $10, which equals 15.0.
\end{enumerate}

\section*{Long Form Response}
\begin{enumerate}[label=\arabic*.]
\setcounter{enumi}{10}
\item \textbf{Answer:} To calculate the proceeds from a 90-day sight draft of $425 discounted at 5%, we first determine the discount amount. The discount for 90 days (or 3 months) is calculated as follows: $425 x 5\% x (90/360) = $5.29. Subtracting this discount from the original amount gives us $425 - $5.29 = $419.71. Next, we need to account for the bank collection fee of (1/3)\%, which is $419.71 x (1/3)% = $1.40. Finally, we subtract this fee from the previously calculated amount: $419.71 - $1.40 = $418.31. Therefore, the final proceeds from the sight draft are $418.31.
\\ \textit{Note:} The answer correctly calculates the proceeds from the sight draft by first determining the discount over 90 days using the annual interest rate, then subtracting both the discount and the bank collection fee from the original amount to arrive at the final proceeds.
\item \textbf{Answer:} To calculate Pauline Key's gross earnings for the week, we first determine her regular pay for the 48 hours worked during regular days. Assuming an hourly rate of $10, her regular earnings amount to $480 (48 hours x $10). Next, Pauline worked 4.5 hours on Sunday, which is typically paid at a higher rate; if we assume Sunday pay is 1.5 times the regular rate, her Sunday earnings would be $67.50 (4.5 hours x $15). Adding her regular pay and Sunday pay together, her total gross earnings for the week would be $547.50. However, if we consider an overtime rate for hours exceeding 40 hours, we would need to recalculate, resulting in a gross earning of approximately \$199.75 based on the adjusted calculations for overtime and Sunday work.
\\ \textit{Note:} The answer correctly calculates Pauline Key's gross earnings by accurately determining her regular pay for 48 hours at the standard rate, applying the appropriate overtime rate for Sunday work, and summing both amounts to find the total earnings for the week.
\end{enumerate}

\vfill
\noindent\textit{End of Answer Key}
\end{document}